%!TEX TS-program = lualatex
%!TEX encoding = UTF-8 Unicode

\documentclass[10pt, a4paper]{awesome-cv}

\colorlet{awesome}{awesome-red}
\setbool{acvSectionColorHighlight}{true}
\renewcommand{\acvHeaderSocialSep}{\quad\textbar\quad}

\usepackage[spanish]{babel}
\usepackage{needspace}
\usepackage{enumitem}

\name{Jeremías}{Likerman}
\position{Geólogo estructural{\enskip\cdotp\enskip} Fracturas naturales{\enskip\cdotp\enskip} Carbonatos y reservorios}
\address{IDEAN, CONICET - Universidad de Buenos Aires}
\mobile{(+54) 230-455-7758}
\email{jlikerman@uba.ar}
\googlescholar{HGA3sQQAAAAJ}{Google Scholar}

\makeatletter
\renewcommand*{\sectionstyle}[1]{{\fontsize{13pt}{15pt}\bodyfont\bfseries\color{text}\@sectioncolor #1}}
\renewcommand*{\subsectionstyle}[1]{{\fontsize{9.4pt}{11pt}\bodyfont\bfseries\scshape\textcolor{graytext}{#1}}}
\renewcommand*{\paragraphstyle}{\fontsize{8.65pt}{11.2pt}\bodyfontlight\upshape\color{text}}
\renewcommand*{\entrytitlestyle}[1]{{\fontsize{9.1pt}{10.8pt}\bodyfont\bfseries\color{darktext} #1}}
\renewcommand*{\entrypositionstyle}[1]{{\fontsize{7.8pt}{9.2pt}\bodyfont\scshape\color{graytext} #1}}
\renewcommand*{\entrydatestyle}[1]{{\fontsize{7.8pt}{9.2pt}\bodyfontlight\slshape\color{graytext} #1}}
\renewcommand*{\entrylocationstyle}[1]{{\fontsize{8pt}{9.5pt}\bodyfontlight\slshape\color{awesome} #1}}
\renewcommand*{\descriptionstyle}[1]{{\fontsize{8.45pt}{10.8pt}\bodyfontlight\upshape\color{text} #1}}
\makeatother

\setlist[itemize]{leftmargin=*,topsep=0.5mm,itemsep=0.15mm,parsep=0mm}

\begin{document}

\makecvheader[C]
\vspace{-5mm}

\cvsection{Perfil orientado al proyecto}
\begin{cvparagraph}
Geólogo estructural especializado en caracterización y predicción de fracturas naturales en sistemas plegados y reservorios energéticos, con más de 15 años de experiencia en campo, análisis estructural, geomecánica y modelado numérico. Su experiencia central en carbonatos naturalmente fracturados incluye la Formación Yacoraite en Tres Cruces, desde sus tesis de licenciatura y doctorado hasta trabajos posteriores sobre redes de fracturas naturales. Integra relevamiento de superficie, análisis multiescala de fracturas, interpretación estructural y criterios predictivos aplicables a reservorios carbonáticos y no convencionales.
\end{cvparagraph}

\vspace{-2mm}
\cvsection{Experiencia relevante en carbonatos, fracturas y reservorios}

\begin{cventries}

\cventry
  {Formación Yacoraite, subcuenca de Tres Cruces}
  {Carbonatos naturalmente fracturados}
  {Jujuy, Argentina}
  {2009 - presente}
  {
    \begin{cvitems}
      \item {Tesis de licenciatura sobre el anticlinal compuesto Tres Cruces y tesis doctoral sobre predicción del fracturamiento en superficies plegadas, con aplicación directa a geometrías de Tres Cruces/Yacoraite.}
      \item {Análisis de redes de fracturas naturales en carbonatos de Yacoraite junto con Clara Correa Luna y Daniel Yagupsky: geometría, orientación, conectividad y relación con estructuras de primer orden.}
    \end{cvitems}
  }

\cventry
  {CONICET - Universidad de Buenos Aires}
  {Investigador en geología estructural y geomecánica}
  {Argentina}
  {2017 - presente}
  {
    \begin{cvitems}
      \item {Desarrollo de modelos predictivos de fracturación natural en pliegues, fallas y superficies geológicas complejas.}
      \item {Integración de observaciones de campo, análisis espacial de fracturas y modelado numérico 2D/3D para evaluar controles estructurales sobre distribución, conectividad y circulación de fluidos.}
    \end{cvitems}
  }

\cventry
  {Vaca Muerta y sistemas análogos}
  {Fracturas naturales en reservorios no convencionales}
  {Cuenca Neuquina, Argentina}
  {2020 - presente}
  {
    \begin{cvitems}
      \item {Coautor de trabajos sobre caracterización de fracturas naturales y conectividad de redes de fracturas en la Formación Vaca Muerta y unidades asociadas.}
      \item {Participación en análisis multiescala de fracturas naturales para Total Austral S.A.; docente del curso profesional \emph{Geología de reservorios no convencionales} e instructor de campo en Vaca Muerta.}
    \end{cvitems}
  }

\cventry
  {Informes técnicos para industria}
  {Caracterización estructural y tectónica aplicada}
  {Argentina}
  {2010 - 2020}
  {
    \begin{cvitems}
      \item {Participación en estudios estructurales y tectónicos para PAE, Petrobras Argentina, Chevron Argentina y Total Austral, con foco en fallas, pliegues, fracturación natural y evolución tectónica aplicada a reservorios.}
    \end{cvitems}
  }

\end{cventries}

\vspace{-3mm}
\cvsection{Formación}
\descriptionstyle{\textbf{Doctor en Ciencias Geológicas, Universidad de Buenos Aires} (2010 - 2015). Tesis: \emph{Predicción del fracturamiento en superficies geológicas plegadas mediante modelado numérico}. Director: Dr. Ernesto Cristallini.}\par
\vspace{0.4mm}
\descriptionstyle{\textbf{Licenciado en Ciencias Geológicas, Universidad de Buenos Aires} (2005 - 2010). Tesis: \emph{Estructura y geología del anticlinal compuesto Tres Cruces, Jujuy, Argentina}. Director: Dr. Ernesto Cristallini.}

\vspace{-3mm}
\cvsection{Competencias aplicadas}
\descriptionstyle{\textbf{Carbonatos y reservorios fracturados:} caracterización de fracturas naturales, conectividad, controles estructurales, integración superficie-subsuelo y criterios predictivos.}\par
\vspace{0.4mm}
\descriptionstyle{\textbf{Geología estructural y campo:} más de 15 años en relevamiento geológico, análisis de pliegues y fallas, cartografía estructural y liderazgo de campañas en los Andes.}\par
\vspace{0.4mm}
\descriptionstyle{\textbf{Modelado y análisis cuantitativo:} modelado numérico y análogo, geomecánica, elementos finitos, análisis espacial/estadístico, Python, R, MATLAB, herramientas geoespaciales y VANT.}\par
\vspace{0.4mm}
\descriptionstyle{\textbf{Idiomas:} español nativo, inglés fluido (C2), italiano (A2), catalán (A2).}

\newpage

\cvsection{Trabajos seleccionados relacionados}
\begin{cvparagraph}
\textbf{Correa-Luna, C., Yagupsky, D. L., Likerman, J. y Barcelona, H.} 2025. \emph{Impact of large-scale structures on fracture network connectivity: Insights into the Vaca Muerta unconventional play, Neuquén Basin, Argentina}. Tectonophysics.

\textbf{Correa-Luna, C., Yagupsky, D. L., Likerman, J. y Barcelona, H.} 2022. \emph{Natural fracture characterization of the Los Catutos Member (Vaca Muerta Formation) in the southern sector of the Sierra de la Vaca Muerta, Neuquén Basin, Argentina}. Journal of South American Earth Sciences.

\textbf{Luna, C. C., Yagupsky, D. L. y Likerman, J.} 2019. \emph{Fracture network analysis of Yacoraite Formation in the Tres Cruces sub-basin, northwestern Argentina}. Journal of South American Earth Sciences.

\textbf{Plotek, B., Likerman, J. y Cristallini, E.} 2024. \emph{Geomechanical modeling of fault-propagation folds: A comparative analysis of finite-element and the trishear kinematic model}. Journal of Structural Geology.

\textbf{Plotek, B., Heckenbach, E., Brune, S., Cristallini, E. y Likerman, J.} 2022. \emph{Kinematics of fault-propagation folding: Analysis of velocity fields in numerical modeling simulations}. Journal of Structural Geology.

\textbf{Likerman, J. y Cristallini, E. O.} 2012. \emph{Predicción de la probabilidad de fracturamiento utilizando métodos estáticos sobre superficies geológicas irregulares}. XV Reunión de Tectónica.

\textbf{Likerman, J., Burlando, J. F., Cristallini, E. O. y Ghiglione, M. C.} 2011. \emph{Estudio de fracturas asociado al sistema de plegamiento en la localidad de Tres Cruces, Provincia de Jujuy}. XVIII Congreso Geológico Argentino.
\end{cvparagraph}

\vspace{-2mm}
\cvsection{Dirección y colaboración científica relevante}
\begin{cventries}
\cventry
  {Tesis doctoral}
  {Fracturación natural en la Formación Vaca Muerta: caracterización y modelado}
  {Correa Luna, Clara}
  {2023}
  {
    \begin{cvitems}
      \item {Dirección: Daniel Yagupsky y Jeremías Likerman. Trabajo enfocado en caracterización y modelado de fracturas naturales en reservorios no convencionales.}
    \end{cvitems}
  }
\cventry
  {Tesis de licenciatura}
  {Geología estructural y fracturación en la cuenca de Tres Cruces}
  {Correa Luna, Clara}
  {2015}
  {
    \begin{cvitems}
      \item {Dirección: Daniel Yagupsky y Jeremías Likerman. Trabajo vinculado a fracturación natural en el sistema estructural de Tres Cruces/Yacoraite.}
    \end{cvitems}
  }
\end{cventries}

\vspace{-2mm}
\cvsection{Cargos actuales y docencia}
\begin{cventries}
\cventry
  {CONICET}
  {Investigador adjunto}
  {Argentina}
  {2017 - presente}
  {
    \begin{cvitems}
      \item {Líneas aplicadas: predicción de fracturas, modelado geomecánico, modelado geodinámico y evaluación de sistemas energéticos.}
    \end{cvitems}
  }
\cventry
  {Departamento de Ciencias Geológicas, Universidad de Buenos Aires}
  {Profesor adjunto}
  {Argentina}
  {2026 - presente}
  {
    \begin{cvitems}
      \item {Asignaturas: Introducción a la Geología, Levantamiento Geológico y Geoestadística.}
    \end{cvitems}
  }
\end{cventries}

\vfill
\begin{center}
{\fontsize{7.8pt}{9pt}\bodyfontlight\color{graytext} CV resumido preparado para proyecto de carbonatos naturalmente fracturados. Mayo de 2026.}
\end{center}

\end{document}
